\documentclass{beamer}

% Theme choice
\usetheme{Madrid} % You can change to e.g., Warsaw, Berlin, CambridgeUS, etc.

% Encoding and font
\usepackage[utf8]{inputenc}
\usepackage[T1]{fontenc}

% Graphics and tables
\usepackage{graphicx}
\usepackage{booktabs}

% Code listings
\usepackage{listings}
\lstset{
basicstyle=\ttfamily\small,
keywordstyle=\color{blue},
commentstyle=\color{gray},
stringstyle=\color{red},
breaklines=true,
frame=single
}

% Math packages
\usepackage{amsmath}
\usepackage{amssymb}

% Colors
\usepackage{xcolor}

% TikZ and PGFPlots
\usepackage{tikz}
\usepackage{pgfplots}
\pgfplotsset{compat=1.18}
\usetikzlibrary{positioning}

% Hyperlinks
\usepackage{hyperref}

% Title information
\title{Chapter 5: Data Structures: Lists and Tuples}
\author{Your Name}
\institute{Your Institution}
\date{\today}

\begin{document}

\frame{\titlepage}

\begin{frame}[fragile]
    \titlepage
\end{frame}

\begin{frame}[fragile]
    \frametitle{Overview of Data Structures in Python}
    \begin{itemize}
        \item Data structures are essential components in programming.
        \item They allow us to organize, manage, and store data efficiently.
        \item In Python, two commonly used data structures are \textbf{Lists} and \textbf{Tuples}.
    \end{itemize}
\end{frame}

\begin{frame}[fragile]
    \frametitle{Key Concepts}
    \begin{enumerate}
        \item \textbf{What are Data Structures?}
            \begin{itemize}
                \item A collection of data values, their relationships, and operations.
                \item Facilitate effective management and processing of large datasets.
            \end{itemize}
        \item \textbf{Importance of Lists and Tuples}
            \begin{itemize}
                \item Group multiple items into a single entity.
                \item Choosing the right structure impacts usability and performance.
            \end{itemize}
    \end{enumerate}
\end{frame}

\begin{frame}[fragile]
    \frametitle{Lists}
    \begin{block}{Definition}
        A list is a mutable sequence of elements that can hold different data types.
    \end{block}
    \begin{itemize}
        \item \textbf{Characteristics:}
            \begin{itemize}
                \item \textbf{Ordered:} Maintains insertion order.
                \item \textbf{Mutable:} Can be modified (add/remove/change elements).
                \item \textbf{Indexing:} Access via 0-based index.
            \end{itemize}
        \item \textbf{Example of a List:}
            \begin{lstlisting}[language=Python]
my_list = [10, "Apple", 3.14, True]
print(my_list[1])  # Output: Apple
            \end{lstlisting}
    \end{itemize}
\end{frame}

\begin{frame}[fragile]
    \frametitle{Tuples}
    \begin{block}{Definition}
        A tuple is an immutable sequence of elements that can also hold multiple data types.
    \end{block}
    \begin{itemize}
        \item \textbf{Characteristics:}
            \begin{itemize}
                \item \textbf{Ordered:} Maintains a definite order.
                \item \textbf{Immutable:} Cannot be changed after creation.
                \item \textbf{Indexing:} Access via 0-based index.
            \end{itemize}
        \item \textbf{Example of a Tuple:}
            \begin{lstlisting}[language=Python]
my_tuple = (20, "Banana", 2.71, False)
print(my_tuple[2])  # Output: 2.71
            \end{lstlisting}
    \end{itemize}
\end{frame}

\begin{frame}[fragile]
    \frametitle{Key Points to Emphasize}
    \begin{itemize}
        \item \textbf{Mutability:}
            \begin{itemize}
                \item Lists can be changed.
                \item Tuples are fixed.
            \end{itemize}
        \item \textbf{Memory Considerations:} 
            \begin{itemize}
                \item Tuples are faster and use less memory because of immutability.
            \end{itemize}
        \item \textbf{Use Cases:}
            \begin{itemize}
                \item Use lists for flexible collections.
                \item Use tuples for constant data.
            \end{itemize}
    \end{itemize}
\end{frame}

\begin{frame}[fragile]
    \frametitle{Conclusion}
    \begin{itemize}
        \item Understanding lists and tuples is fundamental to mastering data handling in Python.
        \item Next, we will explore the specific characteristics of \textbf{Lists} in detail.
        \item Ready to dive deeper?
    \end{itemize}
\end{frame}

\begin{frame}[fragile]
    \frametitle{What are Lists? - Definition}
    A \textbf{list} in Python is an ordered, mutable collection of items that can hold a variety of data types, including integers, strings, and even other lists. 
    Lists allow for the storage and manipulation of multiple items as a single entity.
\end{frame}

\begin{frame}[fragile]
    \frametitle{What are Lists? - Characteristics}
    \begin{enumerate}
        \item \textbf{Ordered}: The items in a list maintain their order, allowing access and iteration in the sequence they were added.
        \item \textbf{Mutable}: Lists can be modified after creation, enabling changes to items, addition of new elements, or removal of existing ones.
        \item \textbf{Indexed}: Each item is assigned an index, starting from 0. Negative indexing supports access to items from the end of the list.
    \end{enumerate}
\end{frame}

\begin{frame}[fragile]
    \frametitle{What are Lists? - Example and Key Points}
    \begin{block}{Example}
        \begin{lstlisting}[language=Python]
# Creating a list
fruits = ['apple', 'banana', 'cherry', 'date']

# Accessing elements using indexing
print(fruits[0])  # Output: apple
print(fruits[-1]) # Output: date

# Modifying elements
fruits[1] = 'blueberry' # Change 'banana' to 'blueberry'
print(fruits) # Output: ['apple', 'blueberry', 'cherry', 'date']
        \end{lstlisting}
    \end{block}
    
    \begin{itemize}
        \item Lists can contain different data types in the same collection:
        \begin{lstlisting}[language=Python]
mixed_list = [1, 'hello', 3.14, True]
        \end{lstlisting}
        
        \item Lists can be nested, meaning you can have lists within lists:
        \begin{lstlisting}[language=Python]
nested_list = [1, [2, 3], ['four', 'five']]
        \end{lstlisting}
        
        \item Lists are dynamic in size:
        \begin{lstlisting}[language=Python]
fruits.append('elderberry')  # Adds 'elderberry' to the end of the list
del fruits[0]                 # Removes the first item from the list
        \end{lstlisting}
    \end{itemize}
\end{frame}

\begin{frame}[fragile]
    \frametitle{Creating and Modifying Lists - Overview}
    \begin{block}{Overview of Lists in Python}
        Lists are versatile data structures that allow you to store, manage, and manipulate groups of related items. Key characteristics include:
    \end{block}
    \begin{itemize}
        \item \textbf{Ordered:} The items maintain the order in which they are added.
        \item \textbf{Mutable:} You can change, add, or remove items after the list has been created.
        \item \textbf{Dynamic:} Lists can grow and shrink as needed.
    \end{itemize}
\end{frame}

\begin{frame}[fragile]
    \frametitle{Creating Lists}
    \begin{block}{Creating Lists}
        To create a list, use square brackets \texttt{[]} and separate items with commas. Lists can contain mixed data types.
    \end{block}
    
    \begin{block}{Syntax}
        \begin{lstlisting}[language=Python]
my_list = [item1, item2, item3]
        \end{lstlisting}
    \end{block}

    \begin{block}{Example}
        \begin{lstlisting}[language=Python]
fruits = ['apple', 'banana', 'cherry']
print(fruits)  # Output: ['apple', 'banana', 'cherry']
        \end{lstlisting}
    \end{block}
\end{frame}

\begin{frame}[fragile]
    \frametitle{Modifying Lists}
    \begin{block}{Adding Elements}
        You can modify lists in several ways using built-in list methods:
    \end{block}
    
    \begin{enumerate}
        \item \textbf{Append:} Add an element to the end of the list.
            \begin{lstlisting}[language=Python]
fruits.append('orange')
print(fruits)  # Output: ['apple', 'banana', 'cherry', 'orange']
            \end{lstlisting}
        
        \item \textbf{Insert:} Add an element at a specific index.
            \begin{lstlisting}[language=Python]
fruits.insert(1, 'mango')
print(fruits)  # Output: ['apple', 'mango', 'banana', 'cherry', 'orange']
            \end{lstlisting}
    \end{enumerate}
\end{frame}

\begin{frame}[fragile]
    \frametitle{Removing and Modifying Elements}
    \begin{block}{Removing Elements}
        To remove elements, you can use:
    \end{block}
    
    \begin{enumerate}
        \item \textbf{Remove:} Remove the first occurrence of a specified value.
            \begin{lstlisting}[language=Python]
fruits.remove('banana')
print(fruits)  # Output: ['apple', 'mango', 'cherry', 'orange']
            \end{lstlisting}
        
        \item \textbf{Pop:} Remove and return an element at a specified index.
            \begin{lstlisting}[language=Python]
last_fruit = fruits.pop()
print(last_fruit)  # Output: 'orange'
print(fruits)      # Output: ['apple', 'mango', 'cherry']
            \end{lstlisting}
    \end{enumerate}
    
    \begin{block}{Modifying Elements}
        Change the value of an element by accessing it via its index:
        \begin{lstlisting}[language=Python]
fruits[0] = 'grape'
print(fruits)  # Output: ['grape', 'mango', 'cherry']
        \end{lstlisting}
    \end{block}
\end{frame}

\begin{frame}[fragile]
    \frametitle{Key Points and Conclusion}
    \begin{block}{Key Points to Remember}
        \begin{itemize}
            \item Lists can store a collection of items, making them robust for many applications.
            \item You can add, remove, and change elements using convenient methods like \texttt{append()}, \texttt{insert()}, \texttt{remove()}, and \texttt{pop()}.
            \item Lists are indexed, meaning you can easily access items via their position in the list (starting with index 0).
        \end{itemize}
    \end{block}
    
    \begin{block}{Conclusion}
        Lists are fundamental to data management in Python. By mastering how to create and modify lists, you’ll be prepared to utilize them effectively in your programming tasks!
    \end{block}
\end{frame}

\begin{frame}[fragile]
    \frametitle{Code Snippet Recap}
    \begin{lstlisting}[language=Python]
# Create a list
fruits = ['apple', 'banana', 'cherry']

# Modify the list
fruits.append('orange')
fruits.insert(1, 'mango')
fruits.remove('banana')
fruits[0] = 'grape'
last_fruit = fruits.pop()
    \end{lstlisting}
    
    Understanding these techniques will pave the way for further exploration of more complex list operations in the upcoming slides!
\end{frame}

\begin{frame}
    \frametitle{Common List Operations}
    \begin{block}{Overview}
        Lists in Python are dynamic and versatile data structures that allow for a variety of operations. Understanding these operations is crucial for effective data handling and manipulation. In this slide, we will cover four key list operations:
        \begin{itemize}
            \item Slicing
            \item Iterating
            \item Sorting
            \item Comprehensions
        \end{itemize}
    \end{block}
\end{frame}

\begin{frame}[fragile]
    \frametitle{Slicing}
    \begin{block}{Definition}
        Slicing allows you to access a subset of elements from a list.
    \end{block}
    
    \textbf{Syntax:} \texttt{list[start:end:step]} \\
    \begin{itemize}
        \item \textbf{start}: Index to begin the slice (inclusive).
        \item \textbf{end}: Index to end the slice (exclusive).
        \item \textbf{step}: Interval between each index. Default is 1.
    \end{itemize}

    \textbf{Example:}
    \begin{lstlisting}[language=Python]
my_list = [0, 1, 2, 3, 4, 5]
slice_example = my_list[1:4]  # Output: [1, 2, 3]
    \end{lstlisting}

    \textbf{Key Points:}
    \begin{itemize}
        \item Negative indices can be used to slice from the end.
        \item Omitting start or end means slicing will continue to the ends of the list.
    \end{itemize}
\end{frame}

\begin{frame}[fragile]
    \frametitle{Iterating and Sorting}
    
    \begin{block}{Iterating}
        \begin{itemize}
            \item \textbf{Definition}: Iterating involves going over each element in a list, typically using loops.
            \item \textbf{For Loop Example:}
        \end{itemize}
        
        \begin{lstlisting}[language=Python]
my_list = [10, 20, 30]
for item in my_list:
    print(item)
        \end{lstlisting}
        
        \textbf{Output:}
        \begin{verbatim}
10
20
30
        \end{verbatim}

        \textbf{Key Points:}
        \begin{itemize}
            \item Efficient way to perform operations on each element.
            \item Can be combined with conditional statements for filtering.
        \end{itemize}
    \end{block}
    
    \begin{block}{Sorting}
        \begin{itemize}
            \item \textbf{Definition}: Sorting arranges the elements in a list in a specified order (ascending or descending).
            \item \textbf{Using the \texttt{sort()} method:}
        \end{itemize}
        
        \begin{lstlisting}[language=Python]
my_list = [3, 1, 4, 2]
my_list.sort()  # Original list is now [1, 2, 3, 4]
        \end{lstlisting}

        \begin{itemize}
            \item \textbf{Using the \texttt{sorted()} function:}
        \end{itemize}
        
        \begin{lstlisting}[language=Python]
my_list = [3, 1, 4, 2]
new_list = sorted(my_list)  # Returns a new sorted list: [1, 2, 3, 4]
        \end{lstlisting}

        \textbf{Key Points:}
        \begin{itemize}
            \item Sorting can be done in reverse by passing \texttt{reverse=True}.
            \item Can also sort custom objects by providing a \texttt{key} function.
        \end{itemize}
    \end{block}
\end{frame}

\begin{frame}[fragile]
    \frametitle{List Comprehensions and Summary}

    \begin{block}{List Comprehensions}
        \begin{itemize}
            \item \textbf{Definition}: A concise way to create lists based on existing lists, allowing for transformation and filtering in a single line of code.
            \item \textbf{Example:}
        \end{itemize}
        
        \begin{lstlisting}[language=Python]
squared_numbers = [x**2 for x in range(5)]  # Output: [0, 1, 4, 9, 16]
        \end{lstlisting}

        \textbf{Key Points:}
        \begin{itemize}
            \item Syntax: \texttt{[expression for item in iterable if condition]}
            \item More readable and often more efficient than traditional loops.
        \end{itemize}
    \end{block}

    \begin{block}{Summary}
        In this slide, we've explored essential operations that enhance your ability to manipulate lists in Python:
        \begin{itemize}
            \item Slicing helps access portions of lists.
            \item Iterating allows processing of each element.
            \item Sorting organizes the data.
            \item List comprehensions provide powerful syntax for creating new lists efficiently.
        \end{itemize}
    \end{block}
\end{frame}

\begin{frame}
    \frametitle{What are Tuples?}
    \begin{block}{Definition}
        Tuples are a built-in data structure in Python used to store an ordered collection of items, 
        which can include various data types such as integers, strings, lists, and even other tuples.
    \end{block}
\end{frame}

\begin{frame}[fragile]
    \frametitle{Characteristics of Tuples}
    \begin{enumerate}
        \item \textbf{Immutability:} Tuples, once created, cannot be altered. This 
        property makes them ideal for storing data that should not change.
        \begin{lstlisting}
my_tuple = (1, 2, 3)
# my_tuple[0] = 5  # Raises TypeError
        \end{lstlisting}

        \item \textbf{Ordered Collection:} Tuples maintain the order of elements, allowing 
        for indexing and slicing.
        \begin{lstlisting}
my_tuple = ('apple', 'banana', 'cherry')
print(my_tuple[1])  # Output: banana
        \end{lstlisting}

        \item \textbf{Heterogeneous:} Tuples can hold diverse data types together.
        \begin{lstlisting}
mixed_tuple = (42, 'hello', 3.14, [1, 2, 3])
        \end{lstlisting}

        \item \textbf{Hashable:} Tuples can be used as keys in dictionaries, unlike lists.
    \end{enumerate}
\end{frame}

\begin{frame}[fragile]
    \frametitle{Use Cases for Tuples}
    \begin{itemize}
        \item \textbf{Returning Multiple Values:} Functions can return multiple values packed in a tuple.
        \begin{lstlisting}
def coordinates():
    return (10, 20)

x, y = coordinates()  # Tuple unpacking
        \end{lstlisting}
        
        \item \textbf{Data Integrity:} To ensure a collection of items does not change, use tuples.
        
        \item \textbf{Dictionary Keys:} Use tuples to create composite keys in dictionaries.
    \end{itemize}
    
    \begin{block}{Key Points to Emphasize}
        \begin{itemize}
            \item Immutability is both a strength (data integrity) and a limitation (no modifications).
            \item Tuples can be faster than lists for certain operations due to their immutable nature.
            \item Tuples generally consume less memory compared to lists, making them suitable for large constant data collections.
        \end{itemize}
    \end{block}
\end{frame}

\begin{frame}[fragile]
    \frametitle{Creating and Using Tuples}
    \begin{block}{Overview}
        Tuples are immutable ordered collections in Python. This slide covers:
        \begin{itemize}
            \item Creating tuples
            \item Accessing tuple elements
            \item Tuple unpacking
        \end{itemize}
    \end{block}
\end{frame}

\begin{frame}[fragile]
    \frametitle{Creating Tuples}
    \begin{block}{Syntax}
        To create a tuple, use parentheses \texttt{()} and separate elements with commas.
        \begin{lstlisting}[language=Python]
my_tuple = (1, 2, 3)
        \end{lstlisting}
    \end{block}
    
    \begin{block}{Points to Remember}
        \begin{itemize}
            \item \textbf{Single Element Tuple:} A single element must include a comma.
                \begin{lstlisting}[language=Python]
single_element_tuple = (5,)
                \end{lstlisting}
            \item \textbf{Empty Tuple:} Create an empty tuple using empty parentheses.
                \begin{lstlisting}[language=Python]
empty_tuple = ()
                \end{lstlisting}
        \end{itemize}
    \end{block}
\end{frame}

\begin{frame}[fragile]
    \frametitle{Accessing Tuple Elements and Unpacking}
    \begin{block}{Accessing Elements}
        Tuples are indexed, allowing access to elements using square brackets. The index starts at 0.
        \begin{lstlisting}[language=Python]
first_element = my_tuple[0]  # Accesses the first element (1)
second_element = my_tuple[1]  # Accesses the second element (2)
        \end{lstlisting}
    \end{block}
    
    \begin{block}{Tuple Unpacking}
        Tuple unpacking assigns elements of a tuple to separate variables in one statement, enhancing clarity.
        \begin{lstlisting}[language=Python]
coordinates = (10, 20)
x, y = coordinates  # x = 10, y = 20
        \end{lstlisting}
    \end{block}
\end{frame}

\begin{frame}[fragile]
    \frametitle{Key Points and Conclusion}
    \begin{block}{Key Points}
        \begin{itemize}
            \item Tuples are immutable; contents cannot be changed.
            \item Can contain mixed data types (integers, strings, lists, etc.).
            \item Useful for fixed collections of items.
        \end{itemize}
    \end{block}
    
    \begin{block}{Conclusion}
        Mastering tuple creation, access methods, and unpacking techniques enhances your coding proficiency and helps create cleaner, more efficient code.
    \end{block}
\end{frame}

\begin{frame}
    \frametitle{Tuples vs. Lists}
    \begin{block}{Overview}
        In this slide, we will explore the differences between lists and tuples in Python. 
        Both are used to store collections of data but serve different purposes and have distinct characteristics. 
        Understanding these differences can help you choose the right data structure for your specific use case.
    \end{block}
\end{frame}

\begin{frame}[fragile]
    \frametitle{Tuples vs. Lists - Definition}
    \begin{itemize}
        \item \textbf{List}: A mutable, ordered collection of items that can contain different data types. Lists are defined using square brackets $[]$.
        \begin{block}{Example}
            \begin{lstlisting}
my_list = [1, 2, 'three', 4.0]
            \end{lstlisting}
        \end{block}
        
        \item \textbf{Tuple}: An immutable, ordered collection of items. Once a tuple is created, its contents cannot be changed. Tuples are defined using parentheses $()$.
        \begin{block}{Example}
            \begin{lstlisting}
my_tuple = (1, 2, 'three', 4.0)
            \end{lstlisting}
        \end{block}
    \end{itemize}
\end{frame}

\begin{frame}[fragile]
    \frametitle{Tuples vs. Lists - Key Differences}
    \begin{enumerate}
        \item \textbf{Mutability}
        \begin{itemize}
            \item \textbf{Lists}:
            \begin{itemize}
                \item Can be modified after creation.
                \item \textbf{Example}:
                \begin{lstlisting}
my_list.append(5)        # Adds 5 to the end
my_list[0] = 'one'       # Changes first element
                \end{lstlisting}
            \end{itemize}
            \item \textbf{Tuples}: 
            \begin{itemize}
                \item Cannot be altered after creation.
                \item \textbf{Example}:
                \begin{lstlisting}
my_tuple = my_tuple + (5,)  # Creates a new tuple
                \end{lstlisting}
            \end{itemize}
        \end{itemize}
        
        \item \textbf{Performance}
        \begin{itemize}
            \item \textbf{Lists}: Slightly slower due to memory overhead.
            \item \textbf{Tuples}: Generally faster and more memory efficient.
        \end{itemize}
        
        \item \textbf{Use Cases}
        \begin{itemize}
            \item \textbf{Choose Lists} for dynamic data or frequent updates.
            \item \textbf{Choose Tuples} for fixed collections, hashable types, or coordinates.
        \end{itemize}
    \end{enumerate}
\end{frame}

\begin{frame}[fragile]
    \frametitle{Conclusion}
    Understanding the differences between lists and tuples will allow you to select the appropriate data structure for your applications based on:
    \begin{itemize}
        \item \textbf{Mutability}: Lists are mutable; tuples are immutable.
        \item \textbf{Performance}: Tuples are generally faster and more memory-efficient than lists.
        \item \textbf{Use Cases}: Choose the structure based on whether the data needs to be modified.
    \end{itemize}
\end{frame}

\begin{frame}[fragile]
    \frametitle{Practical Exercises: Lists and Tuples - Introduction}
    This slide presents hands-on exercises that will help you manipulate \texttt{lists} and \texttt{tuples} in Python. 
    Understanding how to work with these data structures will enhance your programming skills and prepare you for real-world data handling challenges.
\end{frame}

\begin{frame}[fragile]
    \frametitle{Practical Exercises: Lists and Tuples - Key Concepts}
    \begin{itemize}
        \item \textbf{Lists}: Mutable, ordered collections of items. 
        \begin{itemize}
            \item Can contain items of different data types
            \item Items can be modified (e.g., added, removed, or changed)
        \end{itemize}
        \item \textbf{Tuples}: Immutable, ordered collections of items.
        \begin{itemize}
            \item Once defined, their elements cannot be modified
            \item Ideal for fixed collections and ensures data integrity
        \end{itemize}
    \end{itemize}
\end{frame}

\begin{frame}[fragile]
    \frametitle{Practical Exercises: Creating and Modifying a List}
    \textbf{Exercise 1: Creating and Modifying a List}\\
    \textbf{Objective}: Create a list of your favorite fruits and modify it.
    
    \begin{enumerate}
        \item Create a list: 
        \begin{lstlisting}[language=Python]
fruits = ['apple', 'banana', 'cherry']
        \end{lstlisting}
        \item Add a new fruit:
        \begin{lstlisting}[language=Python]
fruits.append('orange')
        \end{lstlisting}
        \item Remove a fruit:
        \begin{lstlisting}[language=Python]
fruits.remove('banana')
        \end{lstlisting}
    \end{enumerate}
    
    \textbf{Expected Result}:
    \begin{itemize}
        \item Original list: \texttt{['apple', 'banana', 'cherry']}
        \item Modified list: \texttt{['apple', 'cherry', 'orange']}
    \end{itemize}
\end{frame}

\begin{frame}[fragile]
    \frametitle{Practical Exercises: Accessing Elements in a Tuple}
    \textbf{Exercise 2: Accessing Elements in a Tuple}\\
    \textbf{Objective}: Create a tuple of integers and access elements by index.
    
    \begin{enumerate}
        \item Create a tuple:
        \begin{lstlisting}[language=Python]
numbers = (1, 2, 3, 4, 5)
        \end{lstlisting}
        \item Access the first and third elements:
        \begin{lstlisting}[language=Python]
first_number = numbers[0]  # Output: 1
third_number = numbers[2]   # Output: 3
        \end{lstlisting}
    \end{enumerate}
    
    \textbf{Key Point}: Remember that tuples are accessed similarly to lists but do not allow for modification of their contents.
\end{frame}

\begin{frame}[fragile]
    \frametitle{Practical Exercises: Slicing Lists and Tuples}
    \textbf{Exercise 3: Slicing Lists and Tuples}\\
    \textbf{Objective}: Use slicing to obtain a subset of items.
    
    \begin{enumerate}
        \item Slicing a list:
        \begin{lstlisting}[language=Python]
fruits = ['apple', 'cherry', 'orange', 'grape']
fruit_slice = fruits[1:3]  # Output: ['cherry', 'orange']
        \end{lstlisting}
        \item Slicing a tuple:
        \begin{lstlisting}[language=Python]
numbers = (1, 2, 3, 4, 5)
number_slice = numbers[1:4]  # Output: (2, 3, 4)
        \end{lstlisting}
    \end{enumerate}
\end{frame}

\begin{frame}[fragile]
    \frametitle{Practical Exercises: Iterating through Lists and Tuples}
    \textbf{Exercise 4: Iterating through Lists and Tuples}\\
    \textbf{Objective}: Use a loop to iterate through the contents of a list and a tuple.
    
    \begin{enumerate}
        \item Looping through a list:
        \begin{lstlisting}[language=Python]
for fruit in fruits:
    print(fruit)
        \end{lstlisting}
        \item Looping through a tuple:
        \begin{lstlisting}[language=Python]
for number in numbers:
    print(number)
        \end{lstlisting}
    \end{enumerate}
    
    \textbf{Expected Output}:
    \begin{itemize}
        \item For \texttt{fruits}: outputs each fruit in the list.
        \item For \texttt{numbers}: outputs each number in the tuple.
    \end{itemize}
\end{frame}

\begin{frame}[fragile]
    \frametitle{Practical Exercises: Summary and Key Takeaways}
    \textbf{Summary}:
    \begin{itemize}
        \item \textbf{Lists}: Useful for maintaining a collection of items that can change dynamically.
        \item \textbf{Tuples}: Ideal for fixed collections of items that should remain constant.
        \item Practice manipulating these structures to become comfortable with their operations, which is crucial for effective programming.
    \end{itemize}
    
    \textbf{Key Takeaways}:
    \begin{itemize}
        \item Understand when to use lists vs. tuples based on mutability needs.
        \item Familiarize yourself with core operations like appending, removing, and slicing.
        \item Use iterative techniques to access elements efficiently.
    \end{itemize}
\end{frame}

\begin{frame}[fragile]
    \frametitle{Best Practices for Using Lists and Tuples - Overview}
    \begin{block}{Understanding Lists and Tuples}
        \begin{itemize}
            \item \textbf{Lists}: Mutable sequences for collections that may need modification.
            \item \textbf{Tuples}: Immutable sequences for fixed collections where data integrity is crucial.
        \end{itemize}
    \end{block}
\end{frame}

\begin{frame}[fragile]
    \frametitle{Best Practices for Using Lists}
    \begin{enumerate}
        \item \textbf{Use Descriptive Variable Names}:
            \begin{itemize}
                \item Choose clear names that describe the contents. 
                \item Example: \texttt{student\_scores} instead of \texttt{list1}.
            \end{itemize}

        \item \textbf{List Comprehensions}:
            \begin{itemize}
                \item Use for concise and efficient creation.
                \item Example:
                \begin{lstlisting}[language=Python]
squares = [x**2 for x in range(10)]
                \end{lstlisting}
            \end{itemize}

        \item \textbf{Keep Lists Homogeneous}:
            \begin{itemize}
                \item Consistent data types for easier processing.
                \item Example: \texttt{students = ['Alice', 'Bob', 'Charlie']}.
            \end{itemize}

        \item \textbf{Avoid Changing Size During Iteration}:
            \begin{itemize}
                \item Avoid modifying lists while iterating. Use copies or list comprehensions.
            \end{itemize}

        \item \textbf{Use Built-in Methods}:
            \begin{itemize}
                \item Leverage Python's built-in methods for tasks.
                \item Example:
                \begin{lstlisting}[language=Python]
my_list.sort()
                \end{lstlisting}
            \end{itemize}
    \end{enumerate}
\end{frame}

\begin{frame}[fragile]
    \frametitle{Best Practices for Using Tuples}
    \begin{enumerate}
        \item \textbf{Choose Tuples for Fixed Data}:
            \begin{itemize}
                \item Ideal for constant data collections.
                \item Example: 
                \begin{lstlisting}[language=Python]
dimensions = (1920, 1080)
                \end{lstlisting}
            \end{itemize}

        \item \textbf{Accessing Tuple Elements}:
            \begin{itemize}
                \item Use indexing (starts at 0).
                \item Example:
                \begin{lstlisting}[language=Python]
color = ('red', 'green', 'blue')
print(color[1])  # Output: green
                \end{lstlisting}
            \end{itemize}

        \item \textbf{Use as Dictionary Keys}:
            \begin{itemize}
                \item Immutable tuples can be dictionary keys.
                \item Example:
                \begin{lstlisting}[language=Python]
point = (2, 3)
my_dict = {point: "This is a point"}
                \end{lstlisting}
            \end{itemize}

        \item \textbf{Unpacking}:
            \begin{itemize}
                \item Assign values to multiple variables in one line.
                \item Example:
                \begin{lstlisting}[language=Python]
x, y = (5, 10)
                \end{lstlisting}
            \end{itemize}
    \end{enumerate}
\end{frame}

\begin{frame}[fragile]
    \frametitle{Summary and Q\&A - Key Points Recap}
    \begin{enumerate}
        \item \textbf{Definitions}:
        \begin{itemize}
            \item \textbf{Lists}: Ordered, mutable, and can contain duplicates. Defined using square brackets \texttt{[]}.
            \item \textbf{Tuples}: Ordered, immutable, and can also contain duplicates. Defined using parentheses \texttt{()}.
        \end{itemize}
        \item \textbf{Mutability}:
        \begin{itemize}
            \item Lists can be modified; tuples cannot.
        \end{itemize}
    \end{enumerate}
\end{frame}

\begin{frame}[fragile]
    \frametitle{Summary and Q\&A - Key Points Recap (cont.)}
    \begin{enumerate}[resume]
        \item \textbf{Accessing Elements}:
        \begin{itemize}
            \item \textbf{Indexing}:
            \begin{lstlisting}
            first_item = my_list[0]    # 1
            second_item = my_tuple[1]   # 2
            \end{lstlisting}
            \item \textbf{Slicing}:
            \begin{lstlisting}
            sub_list = my_list[1:3]     # [2, 3]
            sub_tuple = my_tuple[1:3]   # (2, 3)
            \end{lstlisting}
        \end{itemize}
        \item \textbf{Performance}:
        \begin{itemize}
            \item Tuples can be faster due to their immutability.
        \end{itemize}
    \end{enumerate}
\end{frame}

\begin{frame}[fragile]
    \frametitle{Summary and Q\&A - Applications and Discussion}
    \begin{enumerate}[resume]
        \item \textbf{Applications}:
        \begin{itemize}
            \item Lists: Ideal for collections that require modification.
            \item Tuples: Useful for storing fixed data, like coordinates.
        \end{itemize}
        \item \textbf{Common Methods}:
        \begin{itemize}
            \item List Methods: \texttt{append()}, \texttt{extend()}, \texttt{remove()}, \texttt{pop()}, \texttt{sort()}, \texttt{reverse()}.
            \item Tuple Methods: \texttt{count()}, \texttt{index()}.
        \end{itemize}
    \end{enumerate}
    
    \textbf{Discussion Points:}  
    - When to use a list versus a tuple?  
    - Real-world examples of lists and tuples.  
    - Share experiences about coding with lists/tuples.  
\end{frame}


\end{document}